\section{Related work}
\paragraph{Behavioral and multidimensional evaluation.}
Behavioral testing organizes model evaluation around capabilities and controlled transformations rather than a single held-out score. CheckList formalizes minimum-functionality, invariance and directional-expectation tests for NLP systems \citep{checklist}. System1Bench adopts this separation at a finite-decision interface: candidate order, language realization and irrelevant state fields define invariance-style comparisons, while an operative fact change defines a directional test with recomputed executable labels. The relevant unit is a shared state cluster, not every transformed row treated as independent evidence.

HELM broadens evaluation across scenarios and metrics, emphasizing that accuracy alone leaves substantial behavior unmeasured \citep{helm}. Our setting is narrower: finite typed decisions under explicit candidate and rubric contracts. This narrower interface enables exact prediction matching and shared-state perturbations across native decision heads and code-scored language models. It also requires preserving distinctions between dataset annotations, authored policy references, synthetic teachers and programmatic labels. A large count of decisions cannot substitute for an explicit account of what those references establish.

Recent decision-model studies make a stronger comparison necessary. Legal-document evaluation of Jev and language models already shows that equal aggregate scores can mask different persistent decisions, corrections and regressions \citep{legaldecisions}. Computational-social-science text-annotation work evaluates 18 tasks and many language and decision systems with quality, calibration, routing and cost criteria \citep{cssdecisions}. JevAdvBench evaluates Jev's typed interface across vendor-derived scenarios using exact re-runs, single-edit attacks and confidence-gate outcomes \citep{jevadvbench}. Its primary flip metric deliberately avoids treating model-derived consensus labels as independent gold. System1Bench is not the first cross-domain or stability study. Its narrower complementary question is how correctness transitions and representation sensitivity compare \emph{across model--adapter systems} on fixed downstream decisions with source-specific references, including executable rules where correctness is determined by visible facts. Public labels, authored cases and synthetic teachers remain separate evidence tracks rather than interchangeable ground truth.

Rule-governed decision benchmarks also evaluate whether models apply external rules to facts; RGDT-Bench additionally measures the completeness of stated justifications \citep{rgdtbench}. Our generated policies instead provide an executable oracle for paired answer interventions and do not evaluate explanations or natural-world outcomes. These studies motivate treating broad source coverage as an evaluation substrate, not a novelty claim by itself.

\paragraph{Finite-choice presentation.}
Prior work documents option-identifier bias and benchmark sensitivity to choice order and answer-selection method \citep{mcq,benchmarksensitivity}. Candidate-order sensitivity and code bias are therefore established problems. In our fixed single-token LLM adapter, the codebook control separately changes semantic display order and label-to-code assignment, whereas ordinary reversal changes both. It identifies an interaction in that adapter, not a new general token-bias mechanism or an architectural property. Moreover, first-token scoring can diverge from a model's generated textual answer \citep{firsttoken}; our LLM results characterize a constrained decision interface rather than unconstrained generative capability.

\paragraph{Cross-lingual and contextual behavior.}
XNLI supplies cross-lingual inference examples with shared labels, supporting paired evaluation of language realizations \citep{xnli}. We use the source alignment to compare English and Chinese under the existing question wrapper. Translation realization and wrapper language are part of this intervention; the paired result describes the implemented interface. CLINC150 similarly separates in-scope intent classification from explicit out-of-scope prediction \citep{clinc}. That distinction is operationally relevant because a finite choice among known actions does not by itself supply a reliable reject option.

Long-context results depend on the position of the relevant information as well as the nominal context budget. Lost in the Middle demonstrates this dependence in long-context language-model tasks \citep{middle}. Our needle diagnostic retains source length and insertion position, but also separates the finite-choice intent question from a binary damage question attached to the same state. Grouping by the original needle preserves dependence among length and position variants. The diagnostic concerns explicit fact retrieval, rather than complete long-horizon decision competence.

\paragraph{Selective prediction.}
Selective classification evaluates the trade-off between making predictions and rejecting uncertain cases \citep{selective}. A useful finite-decision model should ideally concentrate its errors among low-confidence cases. We plot empirical risk against coverage within each source, accepting tied confidence values together. The common ranking statistic is maximum class probability. Native decision confidence and a language model's softmax over candidate codes need not share calibration semantics, so selective curves characterize observed error ranking under each adapter rather than interchangeable confidence thresholds.

\paragraph{Inference measurement.}
MLPerf Inference distinguishes operating scenarios and requires performance to be considered with quality constraints \citep{mlperf}. PagedAttention evaluates serving efficiency under an explicitly specified workload and system configuration \citep{vllm}; Sarathi-Serve studies the throughput--latency trade-off in LLM inference \citep{sarathi}. We use these principles to declare timing boundaries, workload shapes, hardware, warm-up and replication. Our local experiment measures an uncached state-to-structured-answer call in a resident process. It excludes model loading, network transport and queueing. Hosted Jev calls expose a different boundary and remain a separate track. This separation is necessary for meaningful interpretation of model-adapter timing rather than a mixed comparison of service and local execution.
